\section{Conclusion}

This project models \textit{Tracks} as a graph-based feasibility problem and
implements the resulting formulation in Python. The grid is represented as an
undirected graph, the railway route is described through selected vertices and
edges, and the path structure is enforced by degree constraints together with
a single-commodity flow formulation. This combination is essential: local
degree rules prevent branches, while the flow certificate eliminates
disconnected cycles and extra components.

The implementation follows the mathematical formulation closely. Instance
files are parsed into normalized data objects, graph helpers construct the
sets used by the MILP, the PuLP/CBC backend solves the feasibility model, and
an independent validator checks the returned solution. Additional tools make
the project compatible with the course workflow by supporting manual test
instances, dataset generation, batch solving, result files, and LaTeX result
tables. The Pygame interface provides a visual layer for inspecting solved
instances and for experimenting with playable boards.

The computational experiments show that the solver behaves correctly on a
controlled generated benchmark from \(5\times5\) to \(16\times16\). All
generated instances in the benchmark are solved to optimality and validated,
while the average solving time increases with the board size and with the
difficulty profile. This confirms both the practical usefulness of the exact
model and the expected growth in difficulty as the graph and route structure
become more complex.

Several extensions remain possible. The generated benchmark could be expanded
to larger grids and more difficulty profiles. The visual interface could also
be improved with richer interaction tools and better puzzle-selection
features. These extensions would make it possible to study the practical
limits of the exact flow formulation while keeping the project focused on the
Tracks model actually required here.
